\section{The Topological Base and Modal Logic Framework}

To formally establish a UTQC, the underlying topological manifold must be rigorously defined. The UOR Atlas achieves this through a categorical mapping of state transitions into a non-pointed $S4$ modal logic structure, utilizing a predefined set of 24 invariant symmetry classes.

\subsection{Formalization of the 24 Symmetry Classes}
The conceptual model (\texttt{docs/conceptual-model}) establishes the exact specifications of the 24 symmetry classes ($\Omega = \{\omega_1, \dots, \omega_{24}\}$). These classes serve as the computational analogues to non-abelian anyonic excitations. By encoding these classes within an $S4$ topological space---characterized by reflexivity and transitivity (idempotence) without requiring point-set specificity---the framework defines a finite Alexandroff topology. The morphisms between these classes constitute the fundamental logic gates of the system.

\subsection{Parametric Proof of Mac Lane Coherence}
A central requirement of any TQC architecture is adherence to the Mac Lane coherence theorem, ensuring that the evaluation of tensor products is topologically invariant. The Pentagon identity governing associativity:
\begin{equation}
\label{eq:pentagon}
    (\alpha_{A,B,C} \otimes id_D) \circ \alpha_{A,B\otimes C,D} \circ (id_A \otimes \alpha_{B,C,D}) = \alpha_{A\otimes B,C,D} \circ \alpha_{A,B,C\otimes D}
\end{equation}
and the Hexagon identity governing commutativity (braiding):
\begin{equation}
\label{eq:hexagon}
    \alpha_{B,C,A} \circ \gamma_{A,B\otimes C} \circ \alpha_{A,B,C} = (\gamma_{A,B} \otimes id_C) \circ \alpha_{B,A,C} \circ (id_B \otimes \gamma_{A,C})
\end{equation}
must be strictly satisfied. 

In theoretical papers, these axioms are often assumed to hold for a proposed physical substrate. In this package, the proof is computational. The validation suite directly asserts Equations (\ref{eq:pentagon}) and (\ref{eq:hexagon}) as executable oracles across all permissible morphisms of the 24 classes. The repository's Continuous Integration (CI) gate executes these algebraic proofs over the combinatorial space, formally proving that the implemented operations form a valid braided monoidal category.
